\documentclass[a4paper,12pt]{report}

\usepackage{amsmath}
\usepackage{subcaption}
\usepackage{hyperref}
\usepackage{xurl}
\usepackage{array}
\usepackage{graphicx}
\usepackage{minted}
\usepackage{multirow}
\usepackage{tikz}
\usepackage{setspace}
\usetikzlibrary{shapes.geometric, arrows}

\doublespacing

\title{How Has the Development of Digital Cryptography Both Influenced and Been
Influenced by Moral and Legal Discussions Surrounding the Right to Privacy?}
\author{Alex Ydens}

% LATEX ARTICLES:
% - https://latex-tutorial.com/tutorials/
% - https://www.overleaf.com/learn/latex/
% - https://www.overleaf.com/learn/latex/LaTeX_Graphics_using_TikZ%3A_A_Tutorial_for_Beginners_(Part_1)%E2%80%94Basic_Drawing
% LATEX NOTES:
% - \section{Section Name} for a section header
% - \subsection{SubSection Name} for a subsection header
% - \paragraph{Paragraph Caption} for a captioned paragraph
% - \subparagraph{Subparagraph Caption} for a captioned subparagraph
% -  surround inline equations in '$' symbols
% - \begin{equation} \end{equation} for multiline equations
% - \begin{equation*} \end{equation*} for non numbered multiline equations
% - \begin{align} \end{align} for multiline equations aligned at the '&' symbols
% - \begin{align*} \end{align*} for non numbered multiline equations aligned at
% the '&' symbols
% -  '\\' in a multiline equation for a new line
% - \label{ident} to label something, \ref{ident} to reference it
% - \left[ \right] for scaled '[]' brackets. same works for '(){}'
% - \begin{matrix} \end{matrix} (with '\\' and '&') for matrices
% - \begin{figure} \end{figure} for a figure
% - \begin{figure}[vals]: '!' for force, 'h' here, 't' top, 'b' bottom, 'p' own
% page. This works for table too, and probably lost of other things.
% - \begin{subfigure}[a-z] \end{subfigure}[a-z] for a subfigure with a label
% (\begin{subfigure}[a] will make figure N(a))
% - \caption{Caption} for a figure caption
% - \label{fig:ident} for a figure label
% - \cite{ident} for a citation
% - \footnote{Footnote} for a footnote
% - \begin{table} \end{table} for a table
% - \url{https://xkcd.com} for a url
% - \textbf{ABC} for bold text
% - \textit{ABC} for italics text
% - \begin{itemize} \end{itemize} with lots of \item for bullet list
% - \begin{enumerate} \end{enumerate} with lots of \item for bullet list

%   \begin{minted}{c}
% #include <stdio.h>
% 
% int main(int argc, char *argv[]) {
%   printf("Hello world!\n");
%   return 0;
% }
% 
%   \end{minted}
% \inputminted{c}{../src/main.c}
% \mintinline{c}|printf("Hello %d\n", 1);|

% EPQ NOTE: 5,000 words target (can be slightly over)

\begin{document}
  \frenchspacing
  \pagenumbering{gobble}
  \maketitle
  \newpage
  \tableofcontents

  \newpage
  \pagenumbering{arabic}
  \section{Abstract}

  \newpage
  \section{Introduction}
These days, trust and security over the internet is becoming more and more
essential. From communicating with banks, to managing sensitive information in a
business, to communicating privately over a service such as WhatsApp, two things
are always necessary: being able to verfiy you are who you say you are, and
being able to send a message that cannot be intercepted and understood in
transit. To accomplish this goal, digital cryptography has had to make several
crucial advances over the last 70 years, but it has been far from smooth
sailing. With so much valuable information at stake, and the risks of giving
people absolute secrecy, research into digital cryptography has marked by
frequently changing attitudes, standards and restrictions. The effect of these
external factors on this field of study, and the implications some of the
advances in it have yielded will be the focus of this document.

  \section{Symmetric Encryption}
The first major breakthrough in digital encryption actually came before
computers in their current form existed, in the 1910s, in the form of the Vernam
cipher. The Vernam cipher, patented by Gilbert Vernam in 1919\cite{vernam}, was
one of the earliest examples of a truly unbreakable cipher, when used correctly.

    \subsection{Using the Vernam Cipher}
To send a message using the Vernam cipher, the sender must first ensure their
message uses a binary encoding. On computers, this is trivial - all data is
encoded in binary on a digital computer - but when originally created, this
required using a tape with holes in it to indicate zeros. Then, a key, which
consists of a binary string of a length greater than or equal to that of the
message, must be bitwise-XORed with the message to obtain the ciphertext. On the
teletype machines of Vernam's time, this required laying the two tapes on top of
each other, and reading any position which had a hole through one tape but not
the other as a one, and all other positions as a zero. To get the plaintext from
the ciphertext, the process could be done in again - just perform a
bitwise-XOR on the ciphertext with the key, and the plaintext should emerge. See
the table below for proof this process works as intended.
      \begin{center}
      \begin{table}[h!]
      \begin{tabular}{ |c|c|c|c| }
        \hline
Plaintext & Key & Ciphertext (Plaintext XOR Key) & Plaintext (Ciphertext XOR
Key) \\
        \hline
0 & 0 & 0 & 0 \\
        \hline
0 & 1 & 1 & 0 \\
        \hline
1 & 0 & 1 & 1 \\
        \hline
1 & 1 & 0 & 1 \\
        \hline
      \end{tabular}
      \caption{Verifying the Vernam Cipher works as intended.}
      \end{table}
      \end{center}
    \subsection{Issues with the Vernam Cipher}
Despite being perfect in theory, and working well with the binary encoding used
by digital computers, the Vernam Cipher has some fairly serious practical issues
when used in isolation. First of all, to ensure perfect secrecy, the key must be
at least as long as the message, completely random, and only known to the sender
and the reciever. This means that, for each message sent, a sender and reciever
would first have to generate a key as long as the message and send it privately.
If the option of privately sending a bit pattern as long as the message was
already available, then the sender and the reciever would not have needed to use
encryption in the first place, and as such the Vernam cipher seems to have
delegated the problem of sending secret messages to key distribution, rather
than having solved the issue.

That being said, initially, there was simply no better way when sending private
digital messages, meaning that those who could afford the cost and were able to
handle the logistics (primalriy large organisations), hired people to deliver
keys in groups of several at a time to those they wanted to communicate with,
and then used the cipher with the delivered keys\cite{singh}. It wasn't perfect,
as the keys could be lost in transit and more would be needed every so often,
but it was better than sending messages regarding sensitive business matters in
plaintext, where anyone without much effort could intercept and read them.

  \section{The DES}
    \subsection{Development}
In response to the increasing use of digital encryption, and a desire to have a
standardised way to utilise it in the government, America's National Bureau of
Standards began to work in collaboration with the public to form a standard for
digital cryptography that could be implemented by whoever needed it, simplifying
the process of cross-organisational communication over encrypted channels.

This was formed initially of an algorithm built on similar principles as
Vernam's cipher created at IBM by Horst Feistel, built into a cipher system that
would hold up well against the attacks of the time. This wasn't that much more
advanced as the Vernam and held the same issues as it, but that could be said of
anything on the market at the time.

The standard was then ran by the National Security Agency of America, who made
some changes. Firstly, they modified it to be less vunerable to differential
cryptanalysis, but more vunerable to brute-force methods, with a decreased key
length of 56 bits\cite{schneier}. This decision was motiviated to ensure that
DES could not be broken by most people, but organisations would be able to put a
lot of time and effort into brute forcing it, thus allowing the US government
the ability to break their citizen's encryption as they saw fit, and ensuring
the standard they released could not be used against them be foreign
adversaries. After this, in 1977, it was released.
    \subsection{Criticism}
The decisision to lower the key length did not go unnoticed amongst the US
public, however. For example, Martin Hellman, who would go on to be one of the
leading figures behind asymmetric encryption, initially spotted that the DES was
vunerable to brute force attacks and, not knowing that this was a conscious
decision by his government, raised the issue to the people behind the standards,
only to be met with a response explaining the reasoning behind the decision. He
made his dissapointment about this clear in a later lecture, stating ``What we
had thought was a technical problem turned out to be political. If we wanted to
improve the security of the standard, we would have to treat it as a political
battle by seeking media coverage and congressional hearings - which we did.".

And Hellman et al were not alone in this. Many critics did not see the weakening
of the security of their communication as a move that would help protect the US
public, seeing as foreign threats could easily develop their own standards.
Furthermore, critics argued that the government having the power to read their
secrets encroached on their right to privacy, treating people as guilty until
proven innocent and allowing for mass surveillance of communications.

It was this frustration with the US government's position on encryption,
according to some\cite{schneier}, that lead to much of the innovation that was
to come in the following years. It certainly was an inspiration for Martin
Hellman\cite{hellman}.

  \section{Diffie-Hellman-Merkle Key Exchange and GCHQ's Equivalent}
The issue of key distribution for long keys, and the insecurity of shorter-key
algorithms such as the one used for the DES standard gave rise to many people to
work on solutions to the key exchange problem, where each solution had its fair
share of upsides and downsides.
    \subsection{Merkle's Puzzles}
Quite possibly the first viable solution to this problem was devised by Ralph
Merkle in 1974. The idea was simple: have the recipient share a piece of
information that can't be used to decipher a message but can be used to encipher
one, so that a sender can encipher their message and send it without requiring
any previous secret communcation, and the reciever can then use the information
they used to generate the public information to decipher the message. In this
way, the key distribution problem is far easier to solve, as the information
containing the keys can be sent publicly without fear of interception.

The idea was to follow the following steps:

    \begin{figure}
      \begin{enumerate}
      \end{enumerate}
\item First, the reciever encrypts \(2^N\) messages of the form: ``The message
is X, the key is Y" with an algorithm that can be brute forced by a sender with
some difficulty, such that brute forcing one would be practical, but brute
forcing \(2^N\) would not. For each message, \(X\) and \(Y\) are chosen
randomly, and the key used to encrypt the message is \(Y\). These messages are
shared publicly, for anyone to use to send their message.
\item The sender then takes a random message from those the recierver has sent,
and uses a brute force method to decrypt it. When they have achiever this, they
send back their message encrypted by the key \(Y\) from the message they
decrypted, followed by the value for \(X\) from that message in plaintext.
\item The reciever can now take the message sent, look through the plaintext
forms of the messages they had encrypted and sent, and find the one which
references the same plaintext \(X\) at the end of the sender's message. This
will point them to the correct \(Y\) to decrypt and read the message.
\item Without knowing which of the \(2^N\) messages the sender chose to brute
force to send their message, an interceptor would need to keep brute forcing
them until the correct one is found, which is impractical.
      \caption{Steps required to use Merkle's Puzzles.}
    \end{figure}

With this system in place, a sender can send the information necessary for
anyone to encrypt a message and send it to them freely without worrying that an
interceptor of this will be able to decipher the following message. As such, key
distribution does not have to take place in secret anymore, so keys can be
bigger, and the Vernam cipher can be used, strengthening messages considerably.

There are two main issues with this system. First of all, it is rather
impractical. To ensure brute forcing \(2^N\) messages cannot occur in a
reasonable time frame, the sender needs to ensure that brute forcing one takes a
considerable amount of time, making the process required of the sender
computationally taxing. Secondly, it is possible the interceptor could get
lucky, and the message the sender chose is one of the first the interceptor
chose to brute force, causing the message to be read without much hassle. As
such, Merkle's Puzzles gave the field of cryptography research hope that the key
distribution problem was solvable, but did not yet provide a solution that would
see much use.
    \subsection{Diffie-Hellman(-Merkle) Key Distribution}

  \section{RSA Encryption}

  \section{Restrictions on RSA and PGP's Legal Issues}

  \section{Quantum Computing's Impact on Digital Cryptography}

  \section{Standards for Post-Quantum Cryptography}

  \newpage
  \section{Appendix}
    \begin{appendix}
      \listoffigures
      \listoftables
      \bibliography{epq}
      \nocite{*}
      \bibliographystyle{ieeetr}
    \end{appendix}
\end{document}
